% Pacotes essenciais
\usepackage{graphicx}

% Fontes
\usepackage{fontspec}
\setmainfont{TeX Gyre Pagella}

% Idioma
\usepackage[brazil]{babel}

% Margens
\usepackage[
    left=3cm,
    right=2.5cm,
    top=2.5cm,
    bottom=3cm
]{geometry}

% Espaçamento
\usepackage{setspace}
\onehalfspacing

% Microtipografia
\usepackage{microtype}

% Cabeçalhos e Rodapés
\usepackage{fancyhdr}
\setlength{\headheight}{16pt}
\pagestyle{fancy}
\fancyhf{} % Limpa
\fancyhead[R]{{\leftmark}}
\fancyhead[L]{\textbf{Apostila de Programação C}}
\renewcommand{\headrulewidth}{0.4pt}
\renewcommand{\footrulewidth}{0pt}

% Primeira página de cada capítulo
\fancypagestyle{plain}{
    \fancyhf{}
    \fancyfoot[C]{\thepage}
    \renewcommand{\headrulewidth}{0pt}
}

% Hyperref
\usepackage[hidelinks]{hyperref}

% Pacote de Códigos
\usepackage{listings}
\lstset{
    language=C,
    basicstyle=\ttfamily\small,
    keywordstyle=\color{blue},
    commentstyle=\color{green!60!black},
    stringstyle=\color{red},
    numbers=left,
    numberstyle=\tiny,
    frame=single,
    breaklines=true
}

% Cores
\usepackage{xcolor}
\definecolor{titulo}{RGB}{20,55,120}

% Caixas (tcolorbox)
\usepackage[most]{tcolorbox}

% Definições de Caixas Pedagógicas
\newtcolorbox{observacao}[1][]{
    colback=blue!5!white,
    colframe=blue!75!black,
    title=Observação,
    fonttitle=\bfseries,
    #1
}

\newtcolorbox{importante}[1][]{
    colback=red!5!white,
    colframe=red!75!black,
    title=Importante,
    fonttitle=\bfseries,
    #1
}

\newtcolorbox{curiosidade}[1][]{
    colback=green!5!white,
    colframe=green!50!black,
    title=Curiosidade,
    fonttitle=\bfseries,
    #1
}

\newtcolorbox{errocomum}[1][]{
    colback=orange!5!white,
    colframe=orange!80!black,
    title=Erro Comum,
    fonttitle=\bfseries,
    #1
}

\newtcolorbox{objetivos}[1][]{
    colback=gray!5!white,
    colframe=gray!50!black,
    title=Objetivos de Aprendizagem,
    fonttitle=\bfseries,
    #1
}

\newtcolorbox{resumo}[1][]{
    colback=yellow!5!white,
    colframe=yellow!60!black,
    title=Resumo,
    fonttitle=\bfseries,
    #1
}

\newtcolorbox{exercicio}[1][]{
    colback=purple!5!white,
    colframe=purple!50!black,
    title=Exercício,
    fonttitle=\bfseries,
    #1
}

\newtcolorbox{desafio}[1][]{
    colback=black!5!white,
    colframe=black!80!white,
    title=Desafio,
    fonttitle=\bfseries,
    #1
}

\newtcolorbox{terminal}[1][]{
    colback=black!90!white,
    colframe=black,
    coltext=white,
    fontupper=\ttfamily,
    title=Terminal,
    fonttitle=\bfseries,
    #1
}

% Títulos
\usepackage{titlesec}
\titleformat{\chapter}
{\Huge\bfseries\color{titulo}}
{\thechapter}
{1em}
{}

\setlength{\parskip}{0.4em}
